\startquestion{2}{二}{填写本问的研究目标，以及相对前问新增的条件。}
\subsubsection*{1）模型的准备}
说明本问实际继承了前问哪些有效资料、定义和结果，以及它们的单位、支持范围与必要限制。新增资料应有明确来源，不能只因为题号相邻就假定所有输出直接可用。

\subsubsection*{2）模型的建立}
围绕本问目标，说明模型选择理由、关键数学关系、参数依据和适用条件。由本问研究审核与你共同确定具体方法；此处没有预设算法，也没有替本问完成研究。

\subsubsection*{3）模型的求解与分析}
填写已确认的方法如何实现和核查，再按本问要求组织真实表图与解释。关键数值从正式结果取用；条件、区间、概率或未达状态应有清晰定义。

本问结束时概括已经回答什么、仍缺什么，并说明其对下一问的数据、假设和整体方向有什么影响。实际不需要该段时按论述自然合并，不机械保留模板提示。
\FloatBarrier
